% Stand 18. Juli 2026
Dieses Kapitel ist wie folgt strukturiert:
\begin{itemize}[topsep=2pt,parsep=1pt,itemsep=0.5pt]
\item[$\circ$] Das Kapitel
\ref{KapitelDatenanalyseHintergrundInfoFeinstaub}
stellt die wesentlichen Informationen rund um das Thema Feinstaub
bereit. Hierbei werden die grundlegenden Definitionen eingeführt,
die möglichen Feinstaub-Quellen beschrieben und in kompakter
Form die Bedeutung von  Feinstaub für die menschliche Gesundheit 
erläutert.
\item[$\circ$] Im Kapitel 
\ref{KapitelDatenanalyseDatenHerkunft} wird zunächst die Herkunft 
der für die Analyse herangezogenen Messdaten aufgezeigt und anschließend
\item[$\circ$] im Kapitel 
\ref{KapitelDatenanalyseDatenStruktur} deren struktureller
Aufbau und die deskriptiv statistischen Charakteristika herausgearbeitet.
\item[$\circ$] Ein wesentlicher Schritt jeder Datenanalyse ist die
Datenbereinigung, die im Kapitel
\ref{KapitelDatenanalyseDatenbereinigung}
beschrieben wird.
\end{itemize}

\subsection{Feinstaub: Definition, Entstehung und Bedeutung für die Gesundheit}
\label{KapitelDatenanalyseHintergrundInfoFeinstaub}

Feinstaub bezeichnet feste und flüssige Schwebeteilchen in der Luft,
 die aufgrund ihrer geringen Größe über längere Zeit in der Atmosphäre
verbleiben können.\\
Unter vereinfachenden Annahmen lässt sich für die Sinkgeschwindigkeit 
$v_s$ von kugel\+för\-migen Partikeln mit Durchmesser $d$  aus der 
Stokes-Gleichung eine quadratische Abhän\-gig\+keit $v_s \propto d^2$ herleiten 
(vgl.\,z.\,B.\,\cite[Kap.\,20]{Landau1987}), was für $d = 2,5\mu m$ zu
 Sinkgeschwindigkeiten im Bereich von Zentimetern pro Stunde und 
 dementsprechend langen atmosphärischen Verweildauern führt.
Insbesondere besitzen kleinere Partikel wesentlich längere
Verweildauern in der Atmosphäre und können daher 
über größere Entfernungen transportiert werden.

\begin{figure}[H] 
\centering
\includegraphics[scale = 0.65]
{Datenanalyse_Einführung_Kategorien Feinstaub.jpg} 
\caption{\small{Atemwegs-Gängigkeit von Feinstaub\protect\footnotemark}}
\label{Abb_Datenanalyse_Einführung_Kategorien Feinstaub}
\end{figure} 

\footnotetext{\textit{Bild entnommen aus}: 
Leitfaden zur Norm ISO 16890 für Luftfiltration der 
Firma \textsc{Mann+Hummel}}
Die von Feinstaub ausgehende Gesundheitsgefährdung hängt maßgeblich davon 
ab, wie weit dieser über den Atemweg 
Nase$\rightarrow$Hals$\rightarrow$Lunge$\rightarrow$Blutkreislauf 
vordringen kann. Diese Atemwegs-Gängigkeit wiederum wird von der Partikelgröße 
des Feinstaubs bestimmt.
Generell gilt: je kleiner die Partikel sind, desto 
tiefer können sie in die Atemwege und die Lunge eindringen 
(vgl.\,Abb.\,\ref{Abb_Datenanalyse_Einführung_Kategorien Feinstaub}).\\

Mit Blick auf empirisch gestützte Studien, in denen die von Feinstaub 
ausgehenden Gesundheitsrisiken aufgegriffen wurden, war es zuerst die
US-amerikanische Umweltschutzbehörde \textbf{E}nvironmental \textbf{P}rotection
\textbf{A}gency, welche den Feinstaub in verschiedene Kategorien einteilte:
\begin{itemize}[topsep=2pt,parsep=1pt,itemsep=0.5pt]
\item[$\circ$] Grobstaub: Partikel mit einem Durchmesser $d$
von $d > 10\mu m$;
\item[$\circ$] $\mathrm{PM}_{10}$: Partikel mit einem
 Durchmesser $d$ von $d \leq 10\mu m$;
\item[$\circ$] $\mathrm{PM}_{2.5}$: Partikel mit einem
 Durchmesser $d$ von $d \leq 2,5\mu m$;
\item[$\circ$] $\mathrm{PM}_{0.1}$: Partikel mit einem
 Durchmesser $d$ von $d \leq 0,1\mu m$\footnote{Da Partikel mit einem Durchmesser 
$d \leq 2,5\mu m$ auch die Bedingung 
$d \leq 10\mu m$ erfüllen, ist der $\mathrm{PM}_{10}$-Wert 
stets größer-gleich dem $\mathrm{PM}_{2.5}$-Wert. Analog folgt, dass 
$\mathrm{PM}_{0.1} \leq \mathrm{PM}_{2.5}$ gilt.}.
\end{itemize}
\vspace*{3mm}

Diese Kategorisierung wurde von der Europäischen Union übernommen 
und die Bestimmung der $\mathrm{PM}_{10}$ und $\mathrm{PM}_{2.5}$ 
Konzentrationen in einer Norm der
\textbf{I}nternational \textbf{S}tandarization 
\textbf{O}rganization festgelegt.\\
%Diese ISO-Norm schreibt zur Bestimmung
%der beiden Feinstaub-Fraktionen die Verwendung eines 
%sog.\,Kaskadenimpaktors vor, dessen \glqq größenselektierender Lufteinlass 
%für Partikel mit einen aerodynamischen Durchmesser von 10 (PM10) bzw. 2,5 (PM2,5) 
%Mikrometer ($\mu m$) eine Abscheidewirksamkeit von 50 Prozent 
%aufweist\grqq\ ($\rightsquigarrow$ 
%\href{https://www.umweltbundesamt.de/service/glossar/pm10}
%{Umweltbundeamt Glossar PM10/PM2,5}).\\

Eine detaillierte Einführung in den Aspekt der Entstehung und
der gesundheitlichen Bedeutung von Feinstaub findet sich 
beispielsweise in den Reports der \textbf{W}orld \textbf{H}ealth 
\textbf{O}rganization \cite{WHO_AirQuality2021} oder des 
\textbf{H}ealth \textbf{E}ffects \textbf{I}nstitute \cite{HEI_2025}, auf die 
sich die nachfolgenden Ausführungen stützen.\\
Feinstaub kann einserseits aus natürlichen Vorgängen hervorgehen 
(z.\,B.\,Waldbrände, aufgewirbelten Bodenstaub).
Gerade in städtischen Gebieten resultiert der größte Teil der 
Feinstaubbelastung allerdings aus menschliche Aktivitäten: 
Reifen- und Bremsabrieb, Verbrennungsprozesse in Motoren, Industrieanlagen 
sowie Heizungen.\\
Für die menschliche Gesundheit besitzt Feinstaub eine hohe Bedeutung,
da die Partikel über die Atemluft aufgenommen werden. Gröbere Partikel 
werden häufig bereits in Nase oder Rachen abgefangen, während kleinere 
Partikel bis in die Bronchien und Lungenbläschen gelangen können. 
Ultrafeine Partikel sind sogar in der Lage, die Blutbahn zu erreichen. Langfristig wird Feinstaub mit einem erhöhten Risiko für 
Herz-Kreislauf-Erkrankungen, chronische Lungenerkrankungen und einer 
verminderten Lebenserwartung in Verbindung gebracht.
Der Report \cite[S.\,5]{HEI_2025} schätzt, dass im Jahr 2023
gesamthaft durch Luftverschmutzung mehr Todesfälle verursacht wurden
als durch Tabakkonsum.
 Aus diesem Grund
zählt die Reduzierung von Feinstaubemissionen zu den wichtigen Zielen 
des Umwelt- und Gesundheitsschutzes.\\
     
Das Bayerische Landesamt für Umwelt und das Umwelt-Ressort der 
Stadt München tragen der großen Bedeutung von Feinstaub für den 
Umwelt- und Gesundheitsschutz dahingehend Rechnung, dass 
landesweit mehrere Messstationen  unterhalten und überwacht werden
($\rightsquigarrow$ \href{https://www.lfu.bayern.de/luft/immissionsmessungen/messwerte/index.htm}
{Aktuelle Messwerte der Luftmess­stati­onen}), hiervon auch vier
Messstationen für Feinstaub (Stand: Juni 2026) im Raum München. \\
Die an diesen vier Stationen erhobenen Messwerte sind \textit{nur} 
als stündliche Mittelwerte verfügbar. 
Abb.\,(\ref{Abb_Datenanalyse_Monatsverlauf_042026_Sensoren_LUA_Stacchus})
zeigt beispielhaft den Verlauf der $\mathrm{PM}_{\mathrm{10}}$- und 
$\mathrm{PM}_{\mathrm{2.5}}$-Messungen an der Messstation am Karlsplatz (Stacchus). 

\begin{figure}[H] 
\centering
\includegraphics[scale = 0.55]
{Datenanalyse_Monatsverlauf_042026_Sensoren_LUA_Stacchus.png} 
\caption{\small{Verlauf der Stunden-Mittelwerte der $\mathrm{PM}_{\mathrm{10}}$- und 
$\mathrm{PM}_{\mathrm{2.5}}$-Messwerte im April 2026 an der Messstation
Karlsplatz (Stacchus) und der festgesetzte Jahresgrenzwert 
(\textcolor{red}{\rule[1mm]{3mm}{0.5mm}})}}
\label{Abb_Datenanalyse_Monatsverlauf_042026_Sensoren_LUA_Stacchus}
\end{figure} 

In Abb.\,(\ref{Abb_Datenanalyse_Monatsverlauf_042026_Sensoren_LUA_Stacchus})
ist zusätzlich der per 
\glqq Verordnung über Luftqualitätsstandards 
und Emissions\-höchst\-mengen\grqq\ 
(\textit{kurz}: BImSchV)\footnote{$\rightsquigarrow$
\href{https://www.lfu.bayern.de/luft/immissionsmessungen/doc/grenzwerte.pdf}
{Grenz- und Zielwerte, kritische Werte, Informations- und Alarmschwellen
nach der 39. BImSchV}}
festgesetzte Jahresgrenzwert für $\mathrm{PM}_{2.5}$ 
von $25\mu g/ m^3$ eingezeichnet 
(\textcolor{red}{\rule[1mm]{3mm}{0.5mm}}).\\

Wie in Kap.\,\ref{KapitelTheorieVariogrammSchätzung} dargelegt, 
können aus $n$ Messpunkten maximal $n\cdot(n-1)/2$ paarweise 
verschiedene Distanzen gebildet werden. 
Mithilfe der vier vom Bayerischen Landesamt für Umwelt in München 
unterhaltenen Messstationen ergeben sich daher nur 6 verschiedene
Distanzen. Dies ermöglicht keine aussagekräftige Schätzung der 
Variogramme, so dass in Folge auch keine Kriging-Analyse basierend 
allein auf diesen Daten möglich ist.  

\subsection{Herkunft und Qualität der Daten}
\label{KapitelDatenanalyseDatenHerkunft}
Untersucht werden Feinstaub-Messungen im Großraum München, die als frei verfügbare 
Daten über das globale Sensor-Netzwerk \textsc{SensorCommunity} 
($\rightsquigarrow$ \href{https://sensor.community/de/}{Sensor Community Home})
zur Verfügung gestellt werden.\\
Die auf dem Portal verfügbaren Sensor-Messungen stammen von 
Nutzern aus über 70 Ländern (Stand 2026) ohne kommerzielle Interessen, 
welche nach einer einmaligen Registrierung die mittels ihrer 
\textit{eigenen} Sensoren erhobenen Messdaten auf das Netzwerk 
hochladen können. \\
Eine Validierung der von den Nutzern hochgeladenen Daten 
oder eine Überprüfung der Qualität der verwendeten 
Sensoren findet \textit{nicht} statt.\\

Die auf \textsc{SensorCommunity} verfügbaren Daten decken die beiden Bereiche
\begin{itemize}[topsep=2pt,parsep=1pt,itemsep=0.5pt]
\item[a) ] Wetter-Bedingungen, also Temperatur, Luftdruck, Luftfeuchtigkeit und
\item[b) ] Feinstaub-Belastung
\end{itemize}
ab. Darüber hinaus finden sich einige Messdateien der radioaktiven 
Umgebungsstrahlung und der Lärmentwicklung. \\
Die historischen Daten sind im Archivbereich 
($\rightsquigarrow$ \href{https://archive.sensor.community/}
{Sensor Community Archiv}) in einer  zeitlich geordneten Struktur in
Komma separierten Format als Messdateien abgelegt 
und werden für unterschiedliche Zeitpunkte für die nachfolgenden Analysen
 herangezogen.
 
Um die typische Anzahl der verfügbaren Messungen zu illustrieren und damit 
die nachfolgenden Überlegungen zu motivieren, seien
hier beispielhaft für den 2.\,April 2026 die Mengenverhältnisse aufgeführt. 
Im Archiv der \textsc{SensorCommunity} sind für diesen Tag 
\begin{itemize}[topsep=2pt,parsep=1pt,itemsep=0.5pt]
\item[$\circ$] insgesamt 19796 Messdateien; davon
\item[$\circ$] 156 Dateien mit zugehörigem Sensor im Großtraum München; davon
\item[$\circ$] 86 Dateien, die Feinstaub-Messungen beinhalten,  
\end{itemize}
verfügbar.\\
Der Download der Daten vom Web-Server der \textsc{SensorCommunity} 
gestaltet sich etwas mühsam: aufgrund zahlreicher 
Verbindungsunterbrechungen dauert das Herunterladen eines 
kompletten Tagessatzes an Feinstaubdaten etwa 3-4 Stunden. \\

Als Mess-Sensor für die Feinstaub-Messungen kommt der SDS011 
($\rightsquigarrow$ \href{https://cdn-reichelt.de/documents/datenblatt/X200/SDS011-DATASHEET.pdf}{SDS011 Datenblatt})
zum Einsatz, der für etwa 30\,EUR erhältlich ist (Stand Juni 2026). 
Hierbei handelt es sich um einen Staubsensor, der die Konzentration des
Feinstaubs in der Luft über die Beeinflussung der Lichtstreuung einer
in den Bausatz integrierten Laserdiode detektiert.\\
Der SDS011 Laser-Staubsensor 
kann Partikel mit einem Durchmesser von mehr als 0,3$\mu m$ erkennen und 
ist ausgelegt für die größen-fraktionierende Detektion von 
\begin{itemize}[topsep=2pt,parsep=1pt,itemsep=0.5pt]
\item[$\circ$] Partikeln mit einem Durchmesser $d \leq 2,5\mu m$ und 
\item[$\circ$] Partikeln mit einem Durchmesser $d \leq 10\mu m$.
\end{itemize}
Die Messungen erfolgen in der Einheit $\mu g/m^3$ und haben gemäß
der Herstellerangabe einen relativen Fehler von 15\%. Als mechanisches Bauteil enthält der SDS011 des Weiteren einen Ventilator, der die
Umgebungsluft ansaugt und zur Lasereinheit leitet. \\

Während nach Herstellerangaben der Mess-Sensor SDS011 in der Lage sein sollte, die beiden Feinstaub-Fraktionen 
$\mathrm{PM}_{10}$ und $\mathrm{PM}_{2.5}$ zu messen, 
zeigten sich bei Vergleichsmessungen mit professionellen
Feinstaub-Messgeräten unter realen Bedingungen Diskrepanzen 
(vgl.\,\cite{Clements2017}, \cite{Patwardhan2021}): 
die Messergebnisse des Messsensors SDS011 korrelieren mit 
den professionell ermittelten $\mathrm{PM}_{x}$
Werte, aber unterschätzen diese tendenziell. 
Nach den Untersuchungen in \cite{Laquai2017} fällt diese
Abweichung bei den $\mathrm{PM}_{10}$
Messwerten größer aus als bei $\mathrm{PM}_{2.5}$
 Messdaten.\\ 
Positiv hingegen fiel in diesen Studien die geringe
Variabilität der SDS011 Sensoren auf.\\
Die zuständige Behörde des australischen Bundesstaates 
New South Wales bezeichnet kostengünstige Feinstaub-Sensoren
wie etwa den SDS011 daher als \glqq indicative air quality
instrument\grqq\ (vgl.\,\cite{NSWDepartment2021}).  \\

Da es in dieser Arbeit primär um die \textit{räumliche Analyse} der 
Daten geht und nicht um deren Interpretierbarkeit im Kontext
anderer Studienergebnisse bedeutet diese Messungenauigkeit 
\textit{keine Einschränkung}.\\
Dennoch, um eventuellen
Fehlinterpretationen der im folgenden präsentierten Ergebnisse
mit anderen, gemäß der Norm durchgeführten Feinstaub-Messungen
vorzubeugen, wird im weiteren Verlauf die Notation 
$\mathrm{P}_{10}$ und $\mathrm{P}_{2.5}$ verwendet für Messwerte,
die mit dem $\mathrm{PM}_{10}$ bzw.\,$\mathrm{PM}_{2.5}$ Messwerte
korrelieren. \\

Die in der Arbeit \cite{Patwardhan2021} erwähnte Unterschätzung
der professionell ermittelten $\mathrm{PM}_x$ durch die 
SDS011 Mess-Sensoren kann auch beim Vergleich der Datenreihen 
der vier Messstationen des Landesumweltamtes mit den Messungen 
von in der Nähe befindlichen SDS011 Sensoren festgestellt werden 
(vgl.\,Abb.\,\ref{Abb_Datenanalyse_Vergleich_PM2x5_Messstation_mit_SDS011_Sensoren}). 

\begin{figure}[H] 
\centering
\includegraphics[scale = 0.55]
{Datenanalyse_Vergleich_PM2x5_Messstation_mit_SDS011_Sensoren.png} 
\caption{\small{Vergleich der stündlichen $\mathrm{PM}_{2.5}$ Mittelwerte  
der Messstation an der Landshuter Allee mit den $\mathrm{P}_{2.5}$
 Mittelwerten von SDS011 Sensoren in rund 500m Entfernung}}
\label{Abb_Datenanalyse_Vergleich_PM2x5_Messstation_mit_SDS011_Sensoren}
\end{figure} 

Es ist deutlich zu erkennen, dass die für eine einstündige Zeitspanne 
ermittelten $\mathrm{PM}_{2.5}$ Mittelwerte 
(\textcolor{orange}{\rule[1mm]{3mm}{0.5mm}})
der Messstation des Landesumweltamts recht deutlich und durchgängig 
oberhalb der ebenfalls über eine Stunde gemittelten $\mathrm{P}_{2.5}$
 Messungen von drei SDS011 Sensoren verlaufen, die in in einer Distanz von 
 etwa 500m positioniert sind. \\
  Diese Diskrepanz zeigt sich auch beim Vergleich der übrigen drei vom
 Landesumweltamt betriebenen Messstationen mit den Daten von in der Nähe
 befindlichen SDS011 Sensoren, die durch eine systematische Abweichung der 
 SDS011 Messungen erklärt werden könnte. \\
Erwähnenswert mit Blick auf 
Abb.\,(\ref{Abb_Datenanalyse_Vergleich_PM2x5_Messstation_mit_SDS011_Sensoren}) 
ist die geringe Variabilität der drei SDS011 Sensoren untereinander; alle drei
Zeitreihen liegen sehr dicht beieinander, was aufgrund ihrer räumlichen 
Nähe zu erwarten sein sollte (erstes Gesetz der Geografie; vgl.\,Kap.\,(\ref{KapitelTheorieKernidee})).


\subsection{Struktur der Daten}  
\label{KapitelDatenanalyseDatenStruktur}  
Die Daten sind in den Messdateien im Wide-Format abgelegt, im Einzelnen:
\begin{itemize}[topsep=2pt,parsep=1pt,itemsep=0.5pt]
\item[$\circ$] \texttt{sensor\_id}: ein natürliche Zahl zur eindeutigen
Identifikation des Sensors;
\item[$\circ$] \texttt{sensor\_type}: eine Zeichenfolge zur Spezifikation
 des verwendeten Sensor-Typs; bei den nachfolgend betrachteten 
 Feinstaub-Messungen ist dies ausnahmslos ein Sensor vom Typ SDS011;
\item[$\circ$] \texttt{location}: eine natürliche Zahl die zur Bestimmung 
der Position des Sensors genutzt werden kann; 
\item[$\circ$] \texttt{lat}: die geodätische Breite (Breitengrad, 
engl.\,\textit{latitude}) angegeben als Gleitkommazahl zur Bestimmung des Winkels 
der räumlichen Position gemessen als Winkel zur Äquatorebene; 
\item[$\circ$] \texttt{lon}: die geodätische Länge (Längengrad, 
engl.\,\textit{longitude}) angegeben als Gleitkommazahl zur Bestimmung des Winkels 
der räumlichen Position bezogen auf den Nullmeridian (0°); 
\item[$\circ$] \texttt{timestamp}: eine Zeichenfolge im Format
YYYY-MM-DDThh:mm:ss zur Bestimmung von Jahr (YYYY), Monat (MM), Tag (DD),
Stunde (hh), Minute (mm) und Sekunde (ss) der Messung;
\item[$\circ$] \texttt{P1}: der $\mathrm{P}_{10}$ Messwert 
(indikativer $\mathrm{PM}_{10}$ Messwert) angegeben in $\mu g/m^3$;
\item[$\circ$] \texttt{durP1}: eine Gleitkommazahl, die bei keiner der Messungen 
mit dem Sensor SDS011 belegt ist; 
\item[$\circ$] \texttt{ratioP1}: eine Gleitkommazahl, die bei keiner der Messungen 
mit dem Sensor SDS011 belegt ist; 
\item[$\circ$] \texttt{P2}: der $\mathrm{P}_{2.5}$ Messwert 
(indikativer $\mathrm{PM}_{2.5}$ Messwert) angegeben in $\mu g/m^3$;
\item[$\circ$] \texttt{durP2}: eine Gleitkommazahl, die bei keiner der Messungen 
mit dem Sensor SDS011 belegt ist; 
\item[$\circ$] \texttt{ratioP2}: eine Gleitkommazahl, die bei keiner der Messungen 
mit dem Sensor SDS011 belegt ist; 
\end{itemize}
 Die beiden geodätischen Koordinaten
Breitengrad und Längengrad sind hierbei nicht im Sexagesimalsystem, also in
Grad (°), Minuten (') und Sekunden ('') angegeben, sondern als Dezimalgrad mit 
drei Nachkommastellen. Dies entspricht einer Granularität von 111m
(vgl.\,z.\,B.\,$\rightsquigarrow$ 
\href{https://wiki.openstreetmap.org/wiki/DE:Genauigkeit_von_Koordinaten}
{Wikipedia: Genauigkeit von Koordinaten}).\\

In Tab.\,(\ref{Tab_Statistsiche_Analyse_Sensor_20260402}) ist beispielhaft 
für eine Messreihe vom 2.\,April 2026 die Struktur-Analyse aufgeführt. Hierbei
fällt mit Blick auf die Zähl-Spalte der NA-Werte (3.\,Spalte) auf, dass für 
alle der insgesamt 569 Messungen an diesem Tag die Einträge
\texttt{durP1}, \texttt{ratioP1}, \texttt{durP2}, \texttt{ratioP2} nicht 
belegt sind. Diese vier Größen werden daher nachfolgend \underline{nicht}
weiter betrachtet. \\
Die Spannweite (engl.\,\textit{span}) gibt die obere und untere Grenze 
(gleichsam das kleinste und größte der Quartile) des beobachteten
Intervalls der Messgröße an. In 
Tab.\,(\ref{Tab_Statistsiche_Analyse_Sensor_20260402}) erfolgt die Angabe des Modus (engl.\,\textit{mode})
im Format häufigster Wert und in Klammern dahinter, wie oft dieser Wert beobachtet wurde. Die Anzahl der verschiedenen beobachteten Werte ist in 
der Spalte \glqq \#Classes\grqq\ wiedergegeben. 
Für metrisch skalierte Variablen sind der arithmetische Mittelwert (engl.\,\textit{mean}) und
Median (engl.\,\textit{median}) angegeben.

\begin{table}[H]
\begin{tabular}{cccccccc}
\hline
Column      & Type   &   \#NA &    Mean &   Median & Span & Mode &   \#Classes \\
\hline \hline
\texttt{sensor\_id} & int & 0 & 6117 &  6117 & [6117,6117] & 6117 (569) & 1 \\
\texttt{sensor\_type} & string & 0 &  nan & nan & [nan,nan] & SDS011 (569) & 1 \\
\texttt{location} & int & 0 & 3089 & 3089 & [3089,3089] & 3089 (569) & 1 \\
\texttt{lat} & float & 0 & 48.20 & 48.20 & [48.20,48.20] & 48.20 (569) & 1 \\
\texttt{lon} & float & 0 & 11.51 & 11.51 & [11.51,11.51] & 11.51 (569) & 1 \\
\texttt{timestamp} & string & 0 & nan &   nan & [nan,nan] & 2026-04... (1) & 569 \\
\texttt{P1} & float  & 0 & 10.64 & 10.27 & [5,20] & 9.80 (9) & 237 \\
\texttt{durP1} & float  & 569 & nan & nan & [nan,nan] & nan (569) & 1 \\
\texttt{ratioP1} & float  & 569 &  nan & nan & [nan,nan] & nan (569) & 1 \\
\texttt{P2} & float  & 0 & 5.01 & 4.9  & [3,8] & 4.80 (15) & 120 \\
\texttt{durP2} & float  & 569 & nan & nan & [nan,nan] & nan (569) & 1 \\
\texttt{ratioP2} & float  & 569 & nan & nan & [nan,nan] & nan (569) & 1 \\
\hline
\end{tabular}
\caption{Statistische Struktur-Analyse beispielhaft für die
Datenreihe eines SDS011 Sensor vom 2.\,April 2026. 
\glqq nan\grqq\ (engl.\,\textit{not a number}) bezeichnet
nicht verfügbare Angaben.}
\label{Tab_Statistsiche_Analyse_Sensor_20260402}
\end{table}

Die Häufigkeits-Analyse der $\mathrm{P}_{10}$-Messwerte desselben Sensors
für den 2.\,April 2026 ist in 
Abb.\,(\ref{Abb_Datenstruktur_Häufigkeitsanalyse_P1_Sensor6117}) 
wiedergegeben.

\begin{figure}[H] 
\centering
\includegraphics[scale = 0.55]
{Datenstruktur_P1_Verteilung_Sensor6117.png} 
\caption{\small{Häufigkeits-Analyse der $\mathrm{P}_{10}$ Messwerte des SDS011 Sensors (ID 6117)}}
\label{Abb_Datenstruktur_Häufigkeitsanalyse_P1_Sensor6117}
\end{figure} 
  
Hier bestätigt das erkennbar links-steile Histogramm die aus dem 
Vergleich von Median und arithmetischen Mittelwert der 
Tab.\,(\ref{Tab_Statistsiche_Analyse_Sensor_20260402})
genährte Vermutung einer nicht symmetrischen Verteilung der Messwerte.\\ 
Das Histogramm der $\mathrm{P}_{2.5}$-Messwerte zeichnet für diesen Sensor qualitativ das 
gleiche Bild.
